\section{The count of rational-exponential pairs}\label{sec:counting}

We count the pairs \((P,v)\) of Section~\ref{sec:setting} for all monic
\(P\) of degree \(p\). Whether an equation can have other meromorphic
waves is a separate question (Section~\ref{sec:completeness}); nothing
in this section or the next depends on the answer.

By Section~\ref{sec:systems}, a normalized profile is a polynomial of
degree \(p\) in \(Q=1/(1+e^{-z})\), without constant term because
\(v\to0\) as \(e^z\to0\). Since \(D^jQ=(-1)^j j!Q^{j+1}\) plus lower
powers of \(Q\), the functions \(Q,DQ,\ldots,D^{p-1}Q\) span the
degree-\(p\) polynomials in \(Q\) with zero constant term, and the
leading pole coefficient \eqref{eq:leadingpole} gives the unique form
\begin{equation}\label{eq:normalized-A}
 v_A=s_pA(D)Q,\qquad
 A(\rho)=\rho^d+\alpha_1\rho^{d-1}+\cdots+\alpha_d,\qquad d=p-1,
\end{equation}
with the scale \(s_p\) of \eqref{eq:leadingpole}; the unknowns of the
count are \(\alpha_1,\ldots,\alpha_d\). This is the standard
singular-part operator form: \(Q=1/2+D\log\cosh(z/2)\), so \(v_A\) is a
constant plus a linear differential operator applied to
\(\log\cosh(z/2)\); such forms, including the KS case, are discussed by
\citet{ConteMusette2003}.

\subsection{Discrete convolution and the matching equations}

For polynomials \(F,G\), let \(F\star G\) be the polynomial with
\((F\star G)(n)=\sum_{j=1}^{n-1}F(j)G(n-j)\) for all integers \(n\ge1\).
It exists because, by the power-sum formulas, the right side is a
polynomial in \(n\). Write \(S_A=A\star A\). For monic \(A\) of degree
\(d\),
\begin{equation}\label{eq:Sleading}
 \deg S_A=2d+1,\qquad
 [\rho^{2d+1}]S_A=\frac{(d!)^2}{(2d+1)!}=\frac{2}{s_p};
\end{equation}
the leading coefficient is the beta integral \(\int_0^1t^d(1-t)^d\dd t\).

\begin{proposition}[Polynomial divisibility]\label{prop:divisibility}
The profile \eqref{eq:normalized-A} solves \eqref{eq:main} for some
monic \(P\) if and only if \(A\) divides \(S_A\).
The operator is unique and is
\begin{equation}\label{eq:operator-A}
 P_A=\frac{s_pS_A}{2A}.
\end{equation}
\end{proposition}

\begin{proof}
In the half-plane \(\operatorname{Re}z<0\), the logistic function is a
geometric series in \(e^z\), and applying \(A(D)\) multiplies the
\(n\)th term by \(A(n)\): \(A(D)Q=\sum_{n\ge1}(-1)^{n-1}A(n)e^{nz}\).
Squaring gives the rational-function identity
\begin{equation}\label{eq:discrete-square}
 (A(D)Q)^2=-S_A(D)Q.
\end{equation}
The map \(F\mapsto F(D)Q\) is injective: its vanishing would give
\(F(n)=0\) for every positive integer \(n\). Consequently
\eqref{eq:main} is equivalent to \(2P_AA=s_pS_A\), which has a
polynomial solution \(P_A\) exactly when \(A\) divides \(S_A\). Equation
\eqref{eq:Sleading} shows that \(P_A\) is monic of degree \(p\).
\end{proof}

Divide \(S_A\) by the monic polynomial \(A\):
\begin{equation}\label{eq:remainder}
 S_A=AC_A+R_A,\qquad
 R_A=\sum_{r=0}^{d-1}R_r(\alpha_1,\ldots,\alpha_d)\rho^r.
\end{equation}
The matching equations are \(R_0=\cdots=R_{d-1}=0\): \(d\) polynomial
equations in the \(d\) unknowns \(\alpha_1,\ldots,\alpha_d\). Their
solutions are in bijection with the normalized pairs \((P_A,v_A)\); we
write \(\Xp\) for the set of solutions. For \(p=2\) and \(A=\rho+a\)
there is the single equation \(R_0=-a(a-1)(a+1)/6=0\) of the
Introduction. Multiplicities in \(\Xp\) refer to these equations.

\paragraph{Simple solutions and multiplicity.}
A solution is \emph{simple} if the Jacobian matrix
\((\partial R_r/\partial\alpha_j)\) is invertible there, like a simple
root, \(f(a)=0\), \(f'(a)\ne0\). Every isolated solution, that is, one
that is not a limit of other solutions, has a \emph{multiplicity}: the
number of solutions near it of the system \(R_r=\varepsilon_r\) for
generic small \(\varepsilon\in\CC^d\), as the double root of \(x^2=0\)
becomes the two roots of \(x^2=\varepsilon\). It is one exactly when the
solution is simple, and it equals the dimension over \(\CC\) of the
convergent power series at that point modulo the combinations
\(\sum_rh_rR_r\) with power series \(h_r\). For finitely many solutions,
the sum of their multiplicities, the \emph{count with multiplicity}, is
the dimension of the polynomials modulo the polynomial combinations
\citep[Sections~4.3 and~5.2]{ArnoldGuseinZadeVarchenko1985},
\citep[Chapter~4, \S2]{CoxLittleOShea2005}; it is the number of distinct
pairs exactly when every solution is simple (see Section~\ref{sec:ks}).

\paragraph{How many solutions can there be?}
Before counting, we ask how many solutions such a system can have at
most. The answer is already \(N_p\). Give \(\alpha_j\) weight \(j\) and
\(\rho\) weight one, so that \(A\) is weighted-homogeneous of weight
\(d\): replacing the roots of \(A\) by \(\lambda\) times themselves
multiplies \(\alpha_j\) by \(\lambda^j\). The polynomial \(S_A\) has
weight at most \(2d+1\), and its part of weight exactly \(2d+1\) is the
\emph{continuous} convolution
\begin{equation}\label{eq:continuous}
 H_A(\rho)=\int_0^\rho A(t)A(\rho-t)\dd t,\qquad
 \int_0^\rho t^i(\rho-t)^j\dd t=\frac{i!\,j!}{(i+j+1)!}\,\rho^{i+j+1}.
\end{equation}
Indeed these are the highest weighted terms of the power sums that
define \(S_A\). Monic division preserves the highest weighted part:
\(A\) is weighted-homogeneous and monic, so each step of long division
subtracts a multiple of \(A\) of the same weight as the term being
removed. Consequently the matching equation \(R_r\) has weight at most
\(2d+1-r\), and
\begin{equation}\label{eq:weights}
 \text{the part of weight }2d+1-r\text{ of }R_r
 \text{ is the coefficient of }\rho^r\text{ in }\remA H_A .
\end{equation}
In words, the discrete convolution \(S_A\) is a lower-order perturbation
of the continuous convolution \(H_A\), and far from the origin the
matching equations look like the continuous problem \(A\mid H_A\).
Solutions of the matching equations could escape to infinity only along
solutions of that problem. \emph{Rigidity} \((\mathrm R_d)\) is the
statement that it has none except the trivial one:
\begin{equation}\label{eq:rigidity}
 A\text{ monic},\quad\deg A=d,\quad A\mid H_A
 \quad\Longrightarrow\quad A(\rho)=\rho^d.
\end{equation}

\begin{proposition}[Weighted B\'ezout bound]\label{prop:bezout}
Let \(p\ge2\) and \(d=p-1\). The matching equations are \(d\) equations
of weighted degrees at most \(2d+1,2d,\ldots,d+2\) in \(d\) unknowns of
weights \(1,2,\ldots,d\). Their isolated solutions, counted with
multiplicity, number at most
\[
 \frac{(d+2)(d+3)\cdots(2d+1)}{d!}=\binom{2d+1}{d}
 =\binom{2p-1}{p-1}=N_p .
\]
If rigidity \((\mathrm R_d)\) holds, that is, if the parts of weight
\(2d+1-r\) of the \(R_r\), which by \eqref{eq:weights} are the
coefficients of \(\remA H_A\), have no common zero other than the
origin, then the matching equations have finitely many solutions, and
counted with multiplicity there are exactly \(N_p\) of them. They are
then \(N_p\) distinct solutions exactly when all of them are simple.
\end{proposition}

For \(p=2,3\) the bound is \(3/1=3\) and \((4\cdot5)/(1\cdot2)=10\). The
substitution \(\alpha_j=y_j^j\) turns weights into ordinary degrees and
counts every solution \(d!\) times; the proof, from B\'ezout's inequality
for isolated solutions \citep[Chapter~VIII]{Mondal2021} and B\'ezout's
theorem \citep{CoxLittleOShea2005}, is Theorem~2.1 of the companion
note. The proof of Theorem~\ref{thm:prime} does not use the proposition:
there the labels give the solutions, and rigidity is a by-product. The
proposition is what remains at other orders (Section~\ref{sec:further}).

For waves, rigidity is a statement about rational solutions. The
dictionary is the Laplace transform, restricted to polynomials.

\begin{lemma}[Laplace correspondence]\label{lem:laplace}
Let \(F\mapsto F^\sharp\) be the linear map on polynomials with
\((\rho^k)^\sharp=k!\,z^{-k-1}\). Then \(D(F^\sharp)=(-\rho F)^\sharp\),
and \(F^\sharp G^\sharp=(F\ast G)^\sharp\) for the continuous convolution
\((F\ast G)(\rho)=\int_0^\rho F(t)G(\rho-t)\dd t\); in particular
\((A^\sharp)^2=(H_A)^\sharp\).
Let \(P\) be monic of degree \(p\), and let \(V\) be a nonconstant
rational solution of \eqref{eq:main} whose only pole is at \(z=0\).
Then \(V\) vanishes at infinity, \(P(0)=0\), and
\[
 V=(-1)^{p-1}s_pA^\sharp,\qquad A\mid H_A,\qquad
 P(\rho)=(-1)^p\,\frac{s_p}{2}\,\frac{H_A}{A}(-\rho),
\]
with \(A\) monic of degree \(p-1\). Conversely, every monic \(A\) of
degree \(p-1\) that divides \(H_A\) gives such a solution. The residue
of \(V\) at zero is \((-1)^{p-1}s_pA(0)\).
\end{lemma}

\begin{proof}
Both rules are checked on monomials; the second is the integral in
\eqref{eq:continuous}. By \eqref{eq:rational-system}, \(V=b+f\) with a
constant \(b\) and with \(f\to0\) at infinity. The constant part of
\eqref{eq:main} is \(b\,(P(0)+b/2)=0\), and the slowest decaying term of
\(f\) enters the rest of the equation with the factor \(P(0)+b\), which
must vanish. Hence \(b=P(0)=0\). By \eqref{eq:leadingpole} the pole has
order \(p\) and leading coefficient \(c_\ast=(-1)^{p-1}(p-1)!\,s_p\), so
\(V=cA^\sharp\) with \(c=(-1)^{p-1}s_p\) and \(A\) monic of degree
\(p-1\). By the two rules, \eqref{eq:main} for \(V\) reads
\(\bigl(cP(-\rho)A+c^2H_A/2\bigr)^\sharp=0\), and the map \(\sharp\) is
injective. This gives the divisibility, the formula for \(P\), and the
converse: the \(P\) defined by the formula is monic of degree \(p\),
because \(H_A\) has the leading coefficient \eqref{eq:Sleading} of
\(S_A\). The residue is the coefficient \(cA(0)\) of \(z^{-1}\).
\end{proof}

For \(A=\rho^{p-1}\) the lemma gives \(V=c_\ast z^{-p}\) and \(P=D^p\).
Rigidity \((\mathrm R_{p-1})\) is therefore the statement that, for
every monic \(P\) of degree \(p\), equation \eqref{eq:main} has no other
nonconstant rational solution with a single pole.

\subsection{The prime-index theorem}\label{sec:prime}

\begin{theorem}[Prime-index enumeration]\label{thm:prime}
Let \(\ell=2p-1\) be prime, and put \(d=p-1\).
\begin{enumerate}
\item[\textup{(a)}] The matching equations have exactly
\(\binom{\ell}{d}=N_p\) solutions over \(\CC\), and every solution is
simple. Thus there are exactly \(N_p\) rational-exponential pairs, and
they are distinct.
\item[\textup{(b)}] The coefficients of every solution are algebraic
numbers, and they are \(\ell\)-integral: they have no \(\ell\) in a
denominator and can be reduced modulo \(\ell\), in the way in which
\(1/11\) becomes \(1\) and \(i\) becomes \(2\) or \(3\) modulo \(5\).
Once a consistent way of reducing algebraic numbers is fixed
(Section~\ref{sec:prime-proof}), reduction modulo \(\ell\) is a
bijection between the solutions and the \(d\)-element subsets \(L\) of
the field \(\FF_\ell\) of integers modulo \(\ell\):
\begin{equation}\label{eq:prime-labels}
 \overline A(\rho)=\prod_{s\in L}(\rho-s).
\end{equation}
\item[\textup{(c)}] Rigidity \((\mathrm R_d)\) holds.
\end{enumerate}
\end{theorem}

We call \(L\) the \emph{label} of the pair. It depends on the way of
reducing algebraic numbers, an embedding fixed in the proof, only as
\(i\) is \(2\) or \(3\) modulo \(5\): another embedding gives a solution
the label that the first one gives to a conjugate solution. Whether a
label contains \(0\), and the labels of solutions with rational
coefficients, do not depend on it. The labels describe reductions of
algebraic coefficients; they do not assert that the complex roots of
\(A\), or the characteristic roots of \(P_A\), are integers.

\paragraph{Example: order three.}\label{sec:ks-rational}
Let \(p=3\), \(s_3=60\), and \(A(\rho)=\rho^2+a\rho+b\), so that
\((a,b)=(\alpha_1,\alpha_2)\). The profile and monic operator are
\begin{align}
 \frac{v_A}{60}&=2Q^3-(a+3)Q^2+(a+b+1)Q, \label{eq:ksprofile}\\
 P_A(\rho)&=\rho^3+4a\rho^2+(a^2+19b)\rho-a^3+7ab-5a-30b.
   \label{eq:ksoperator}
\end{align}
The remainder of \(S_A\) is \((\widehat R_1\rho+\widehat R_0)/30\) with
\begin{equation}\label{eq:ksmatching}
 \widehat R_0=30R_0=ab(a^2-7b+5),\qquad
 \widehat R_1=30R_1=a^4-8a^2b+11b^2+10b-1.
\end{equation}
We solve \(\widehat R_0=\widehat R_1=0\) by cases on the factors of
\(\widehat R_0\). If \(b=0\), then \(a^4=1\). If \(a=0\), then
\((b+1)(11b-1)=0\). In the remaining case \(a^2=7b-5\), and the second
equation becomes \(4(b-2)(b-3)=0\). This exhausts the system and gives
Table~\ref{tab:ks}. The determinant
\(\det\partial(\widehat R_0,\widehat R_1)/\partial(a,b)\) is nonzero at
every point, so all \(N_3=10\) pairs are simple, as
Theorem~\ref{thm:prime}(a) states for \(2p-1=5\).

\begin{table}[htbp]
\centering
\caption{All ten rational-exponential pairs at order three. The label
is the set of roots of \(A\) modulo \(5\); for the two complex pairs it
depends on the choice of \(i\) modulo \(5\). For the equation of the
Introduction, \((b_1,b_2,b_3)=(1,\sigma,1)\), the coefficients \(B=4a\)
of \(\rho^2\) and \(\mu=a^2+19b\) of \(\rho\) in \eqref{eq:ksoperator}
are \(\sigma/\kappa\) and \(1/\kappa^2\) by \eqref{eq:coefficients}, so
\(\sigma^2=B^2/\mu=16a^2/(a^2+19b)\). In the rows \((\pm i,0)\) and
\((0,-1)\), \(\mu<0\) and the rate \(\kappa\) is imaginary. The last
column describes \(v_A\) on the real axis, at unit rate.}
\label{tab:ks}
\begin{tabular}{@{}cclccl@{}}
\toprule
\(a\)&\(b\)&Labels&\(|\sigma|\)&Determinant&Real profiles\\
\midrule
\(\pm1\)&\(0\)&\(\{0,4\}\), \(\{0,1\}\)&\(4\)&\(-24\)&Pulses\\
\(\pm i\)&\(0\)&\(\{0,2\}\), \(\{0,3\}\)&\(4\)&\(-16\)&Not real\\
\(0\)&\(-1\)&\(\{1,4\}\)&\(0\)&\(144\)&Front\\
\(0\)&\(1/11\)&\(\{2,3\}\)&\(0\)&\(576/121\)&Front\\
\(\pm3\)&\(2\)&\(\{3,4\}\), \(\{1,2\}\)&\(12/\sqrt{47}\)&\(-144\)&Fronts\\
\(\pm4\)&\(3\)&\(\{2,4\}\), \(\{1,3\}\)&\(16/\sqrt{73}\)&\(384\)&Fronts\\
\bottomrule
\end{tabular}
\end{table}

The ten labels are the ten two-element subsets of \(\{0,1,2,3,4\}\). By
\eqref{eq:ksprofile}, \(v_A/60\to b\) as \(z\to+\infty\), so a pair is a
pulse exactly when \(b=0\). The map \((a,b)\mapsto(-a,b)\) is the mirror
image of the Introduction, and it replaces a label \(L\) by \(-L\)
(Section~\ref{sec:reflection}). Eight of the ten pairs have real
coefficients, and all eight are smooth and bounded on the real axis: two
pulses and six fronts. For example, \(A=(\rho-1)(\rho-2)\) gives
\(v_A=120Q^3\) and \(P_A=(\rho-3)(\rho-4)(\rho-5)\). These are the
classical solutions in hyperbolic functions
\citep{KuramotoTsuzuki1976,Kudryashov1988,Kudryashov1989,Kudryashov1990,Kudryashov1991},
tabulated completely by \citet[Table~1]{KudryashovZargaryan1996} and
\citet[Table~1]{ConteMusette2009}. Their tables contain six cases
because they identify the two signs of the dispersion coefficient: the
four mirror pairs \((\pm a,b)\) and the two pairs with \(a=0\).
Exhaustion of this list follows from \citet{Eremenko2006};
\citet{Demina2020} classifies all irreducible algebraic invariants of
the same profile equation, and later rederivations by ansatz methods
coincide with these waves \citep{KudryashovSoukharev2009}. The pairs
\((\pm1,0)\) and \((\pm i,0)\) belong to the same equations, those with
\(\sigma^2=16\), at the rates \(\kappa=\pm1\) and \(\kappa=\pm i\)
(Section~\ref{sec:fixed}). For a real equation the second wave is thus
periodic in \(x\), and its real form has poles on the real axis.

\subsection{Proof of the prime-index theorem}\label{sec:prime-proof}

The idea was described in the Introduction. A reader who accepts the
theorem can go to Section~\ref{sec:labels}: the rest of the paper uses
only its statement, and the proof of Proposition~\ref{cor:endpoints} uses
\eqref{eq:prime-collapse} and fact \textup{(V)} below once more.

\paragraph{Reduction modulo \(\ell\) of algebraic numbers.}
We use four standard facts. Rational numbers without \(\ell\) in the
denominator form a ring \(\ZZ_{(\ell)}\). Think of \(\ell\) as a small
parameter \(\varepsilon\): the valuation is the order of a quantity in
\(\varepsilon\), and reduction modulo \(\ell\) sets \(\varepsilon=0\).

\emph{\textup{(V)} Valuation and reduction.} Let \(K\) be an algebraic
closure of the field \(\QQ_\ell\) of \(\ell\)-adic numbers. It is
algebraically closed, contains \(\QQ\), and carries the \(\ell\)-adic
valuation \(\operatorname{ord}\colon K\to\QQ\cup\{\infty\}\)
\citep[Chapter~III]{Koblitz1984},
\citep[Chapter~II, Theorem~(4.8)]{Neukirch1999}: the exponent of
\(\ell\) in a nonzero rational number, infinite only at \(0\), with
\begin{equation}\label{eq:valuation-rules}
 \operatorname{ord}(xy)=\operatorname{ord}(x)+\operatorname{ord}(y),\qquad
 \operatorname{ord}(x+y)\ge\min\{\operatorname{ord}(x),\operatorname{ord}(y)\}.
\end{equation}
We fix an embedding of the field \(\overline\QQ\subset\CC\) of algebraic
numbers in \(K\), unique up to an automorphism of \(\overline\QQ\), and
regard algebraic numbers as elements of \(K\). Those \(x\in K\) with
\(\operatorname{ord}(x)\ge0\) are \emph{\(\ell\)-integral}. They form a
ring \(\mathcal O\supset\ZZ_{(\ell)}\), and \emph{reduction modulo
\(\ell\)}, \(x\mapsto\overline x\), maps \(\mathcal O\) onto its residue
field \(\Bbbk\supset\FF_\ell\): it respects sums and products, it is the
usual reduction on \(\ZZ_{(\ell)}\), and \(\overline x=0\) exactly when
\(\operatorname{ord}(x)>0\). Polynomials are reduced coefficient by
coefficient. Division by a monic polynomial uses only sums and products
of coefficients: it commutes with substituting values for the
\(\alpha_j\) and with reduction, and it stays within
\(\ZZ_{(\ell)}[\alpha_1,\ldots,\alpha_d]\) or \(\mathcal O\). A monic
polynomial with \(\ell\)-integral roots has \(\ell\)-integral
coefficients.

\emph{\textup{(H)} Hensel's lemma.} Let \(G=(G_1,\ldots,G_d)\) be
polynomials in \(d\) unknowns with coefficients in \(\ZZ_{(\ell)}\). If
\(\overline a\in\FF_\ell^d\) solves the reduced system and its Jacobian
determinant is not zero there, then exactly one \(a\in\mathcal O^d\) has
\(G(a)=0\) and reduction \(\overline a\). This is Newton's method from
an integer vector that reduces to \(\overline a\), just as a simple root
at \(\varepsilon=0\) continues uniquely to a power series in
\(\varepsilon\): the multivariable Hensel lemma of \citet{ConradHensel},
applied in the finite extensions of \(\QQ_\ell\) in \(K\), which are
complete \citep[Chapter~II, Theorem~(4.8)]{Neukirch1999}.

\emph{\textup{(N)} The weak Nullstellensatz.} Polynomial equations with
rational coefficients that have a solution in \(\CC^d\) have one in
\(\overline\QQ^d\) \citep[Chapter~4, \S1]{CoxLittleOSheaIVA}.

\emph{\textup{(T)} Finite systems.} If \(d\) polynomial equations in
\(d\) unknowns with rational coefficients have only finitely many
solutions in \(K^d\), then these lie in \(\overline\QQ^d\), and they are
also the solutions in \(\CC^d\), with the Jacobian determinant nonzero
at the same ones. Indeed, by finiteness a nonzero polynomial in each
\(x_i\) alone is a combination of the equations
\citep[Chapter~5, \S3]{CoxLittleOSheaIVA}; finding one of bounded degree
is a linear problem over \(\QQ\), so one with rational coefficients
exists, and it confines \(x_i\) to its roots, in \(K\) as in \(\CC\).

\begin{proof}[Proof of Theorem~\ref{thm:prime}]
The proof has four steps: a congruence for \(S_A\) modulo \(\ell\); the
solution of the reduced equations; integrality, a dominant balance
showing that every solution can be reduced, and then reduces to one of
them; and lifting, Newton's method in the form of Hensel's lemma,
showing that each of them comes from exactly one solution. Solutions are
first taken in \(K^d\), with \(K\) as in \textup{(V)}; the last step
returns to \(\CC\).

\emph{Step 1: the congruence.} Induction using \((D-n)Q^n=-nQ^{n+1}\)
gives \((n-1)!\,Q^n=(-1)^{n-1}\prod_{j=1}^{n-1}(D-j)Q\). Write
\((A(D)Q)^2=\sum_{n=2}^{2p}\gamma_nQ^n\). Each \(\gamma_n\) is an
integer polynomial in the coefficients of \(A\), and
\(\gamma_{2p}=(d!)^2\). By \eqref{eq:discrete-square},
\[
 S_A(\rho)=\sum_{n=2}^{2p}\frac{(-1)^n\gamma_n}{(n-1)!}
       \prod_{j=1}^{n-1}(\rho-j).
\]
Only the last denominator, \((2p-1)!=\ell!\), contains \(\ell\).
Consequently \(\ell S_A\) has coefficients in
\(\ZZ_{(\ell)}[\alpha_1,\ldots,\alpha_d]\), and
\begin{equation}\label{eq:prime-collapse}
 \ell S_A(\rho)\equiv\frac{(d!)^2}{(\ell-1)!}\prod_{j=1}^{\ell}(\rho-j)
 \equiv(-1)^d(\rho^\ell-\rho)\pmod\ell.
\end{equation}
The scalar follows by pairing \(j\) with \(\ell-j\) in \((\ell-1)!\).
The product is \(\rho^\ell-\rho\) because, by Fermat's little theorem,
every residue \(s\) modulo \(\ell\) satisfies \(s^\ell=s\). The right
side does not depend on \(A\): this is what makes the prime index
special. Taking the parts of weight \(\ell=2d+1\) on both sides gives
the corresponding congruence for the continuous convolution,
\begin{equation}\label{eq:prime-continuous}
 \ell H_A(\rho)\equiv(-1)^d\rho^\ell\pmod\ell .
\end{equation}
Both congruences are between polynomials in \(\rho\) and the
\(\alpha_j\) with coefficients in \(\ZZ_{(\ell)}\), so they may be
evaluated at any \(\ell\)-integral values of the \(\alpha_j\) and of
\(\rho\), and then reduced.

\emph{Step 2: the reduced equations.} By \textup{(V)} the scaled
remainders \(\ell R_r\) have coefficients in \(\ZZ_{(\ell)}\), and by
\eqref{eq:prime-collapse} their reductions are the coefficients of the
remainder of \((-1)^d(\rho^\ell-\rho)\) on division by \(\overline A\).
The reduced equations therefore say that \(\overline A\) divides
\(\rho^\ell-\rho\). Now \(\rho^\ell-\rho=\prod_{s\in\FF_\ell}(\rho-s)\),
and polynomials over a field factor uniquely. Over any field that
contains \(\FF_\ell\), in particular over \(\Bbbk\), the monic divisors
of degree \(d\) are therefore the \(\binom{\ell}{d}\) polynomials
\eqref{eq:prime-labels}, and their coefficients lie in \(\FF_\ell\).
Each of them is a simple solution of the reduced equations. Write
\(\rho^\ell-\rho=\overline A\,\overline C\). For a variation \(h\) with
\(\deg h<d\), the derivative of the remainder is, up to the nonzero
scalar \((-1)^d\),
\begin{equation}\label{eq:prime-jacobian}
 h\longmapsto-\operatorname{rem}_{\overline A}(\overline Ch),
\end{equation}
because the dividend in \eqref{eq:prime-collapse} does not vary with
\(A\). Since \(\rho^\ell-\rho\) has no repeated factor, \(\overline A\)
and \(\overline C\) have no common factor. If \(\overline A\) divides
\(\overline Ch\), it therefore divides \(h\), and then \(h=0\) because
\(\deg h<d\). The linear map \eqref{eq:prime-jacobian} is thus
injective, hence invertible.

\emph{Step 3: integrality and rigidity.} A solution could have \(\ell\)
in a denominator, and then it could not be reduced. We show that this
does not happen. The argument is a dominant balance, with \(\ell\) in
the role of the small parameter: a solution that is large is rescaled to
size one, and the leading-order problem, which is the continuous one,
then leads to a contradiction. For \(\lambda\ne0\) and a monic
polynomial \(A\) of degree \(d\), put
\(A_\lambda(\rho)=\lambda^{-d}A(\lambda\rho)\). This is the monic
polynomial whose roots are those of \(A\) divided by \(\lambda\), and
its coefficients are \(\alpha'_j=\alpha_j/\lambda^j\). A polynomial
\(T\) in the \(\alpha_j\) and \(\rho\) that is weighted-homogeneous of
weight \(w\) satisfies
\(T(\alpha;\lambda\rho)=\lambda^{w}T(\alpha';\rho)\).

Let \(A\) be a solution of the matching equations with coefficients in
\(K\), and suppose that some root of \(A\) is not \(\ell\)-integral. Let
\(\lambda\) be a root of least valuation, so that
\(\theta=\operatorname{ord}(\lambda)<0\), and put \(B=A_\lambda\). The
roots of \(B\) are \(\ell\)-integral, hence so are its coefficients
\(\alpha'_j\), by \textup{(V)}. Decompose
\(\ell S_A=\sum_{w\le\ell}T^{(w)}\) into its parts of weight \(w\). They
have coefficients in \(\ZZ_{(\ell)}\) by Step~1, and
\(T^{(\ell)}=\ell H_A\). Since \(A\) divides \(S_A\) and
\(A(\lambda)=0\), we have \(S_A(\lambda)=0\). Dividing this equation by
\(\lambda^{\ell}\) gives
\[
 \ell H_B(1)=-\sum_{w<\ell}\lambda^{w-\ell}\,T^{(w)}(\alpha';1).
\]
Each term on the right has positive valuation, because
\(\operatorname{ord}(\lambda^{w-\ell})=(w-\ell)\theta>0\) and
\(T^{(w)}(\alpha';1)\) is \(\ell\)-integral. By
\eqref{eq:valuation-rules} the right side reduces to zero. The left side
reduces to \((-1)^d\ne0\), by \eqref{eq:prime-continuous} evaluated at
\(\alpha=\alpha'\) and \(\rho=1\). This contradiction shows that every
root of \(A\) is \(\ell\)-integral, and so is every coefficient. The
solution can be reduced, and by Step~2 its reduction is one of the
polynomials \eqref{eq:prime-labels}.

Rigidity is proved in the same way. Suppose that \(A\ne\rho^d\) is monic
of degree \(d\), with complex coefficients, and divides \(H_A\). Since
\(H_{A_\lambda}(\rho)=\lambda^{-\ell}H_A(\lambda\rho)\), every
\(A_\lambda\) has the same property, so we may assume that one of the
coefficients \(\alpha_j\) equals \(1\). By \textup{(N)}, applied to the
equations \(\remA H_A=0\) together with this normalization, we may
assume that the coefficients of \(A\) lie in \(\overline\QQ\subset K\).
Let \(\lambda\) be a root of \(A\) of least valuation; it is not zero,
because \(A\ne\rho^d\). As before, \(B=A_\lambda\) has \(\ell\)-integral
coefficients. Since \(A\) divides \(H_A\) and \(A(\lambda)=0\), we have
\(0=\ell H_A(\lambda)=\lambda^{\ell}\,\ell H_B(1)\). But \(\ell H_B(1)\)
reduces to \((-1)^d\ne0\), by \eqref{eq:prime-continuous}. This
contradiction proves~(c).

\emph{Step 4: lifting.} The system \((\ell R_r)_{0\le r<d}\) has
coefficients in \(\ZZ_{(\ell)}\), and by Step~2 each of the
\(\binom{\ell}{d}\) reduced solutions lies in \(\FF_\ell^d\) and is
simple. By Hensel's lemma \textup{(H)}, each of them is the reduction of
exactly one solution in \(\mathcal O^d\). By Step~3 every solution in
\(K^d\) lies in \(\mathcal O^d\), and its reduction is one of the
reduced solutions. Hence reduction modulo \(\ell\) is a bijection
between the solutions in \(K^d\) and the reduced solutions. There are
exactly \(\binom{\ell}{d}\) solutions in \(K^d\), and the Jacobian
determinant is nonzero at each of them, because its reduction is
nonzero. By \textup{(T)}, the matching equations have the same
\(\binom{\ell}{d}\) solutions over \(\CC\), with coordinates in
\(\overline\QQ\), all simple. This proves (a) and (b); part~(c) was
proved in Step~3.
\end{proof}

\begin{remark}[Formal verification]\label{rem:lean}
Parts (a) and (c) of Theorem~\ref{thm:prime} have been verified in
Lean~4 against the Mathlib library, both for the matching equations and
for the pairs \((P,v)\) themselves (Appendix~\ref{app:verification}).
The formal proof follows the route above, and like the proof above it
normalizes the roots of \(A\).
\end{remark}

\subsection{Reading the labels: pulses and fronts}\label{sec:labels}

As \(z\to+\infty\) we have \(Q\to1\) and \(D^jQ\to0\) for \(j\ge1\), so
\(F(D)Q\to F(0)\) for every polynomial \(F\). The limits of \(v_A\) as
\(z\to\mp\infty\) are therefore \(0\) and \(s_pA(0)\): pairs with
\(A(0)=0\) are pulses, and the others are fronts. The pulses are needed
separately, because the elliptic pairs of
Section~\ref{sec:elliptic-counting} degenerate onto them.

Letting \(z\to+\infty\) in \eqref{eq:discrete-square} gives
\(S_A(0)=-A(0)^2\). Since \(S_A(0)=A(0)C_A(0)+R_0\) by
\eqref{eq:remainder}, this gives the polynomial identity
\begin{equation}\label{eq:constant-remainder}
 R_0=-A(0)\bigl(A(0)+C_A(0)\bigr).
\end{equation}
On the hyperplane \(\alpha_d=A(0)=0\) the equation \(R_0=0\) thus holds
identically, and the pulses are the solutions of the \emph{pulse equations}
\begin{equation}\label{eq:pulse-equations}
 R_1=\cdots=R_{d-1}=0\qquad\text{in the unknowns }
 \alpha_1,\ldots,\alpha_{d-1},\quad\alpha_d=0 .
\end{equation}
We write \(\mathcal Z_p\) for the set of their solutions; multiplicities
of points of \(\mathcal Z_p\) refer to the pulse equations.

\begin{proposition}[Pulses and fronts]\label{cor:endpoints}
Let \(\ell=2p-1\) be prime. A pair is a pulse if and only if its label
contains \(0\). Consequently exactly \(S_p:=\binom{2p-2}{p-2}\)
of the \(N_p\) pairs are pulses, and the remaining
\(\binom{2p-2}{p-1}\) are fronts. Every pulse is a simple solution of
the pulse equations \eqref{eq:pulse-equations}, and
\(P_A(0)\ne0\) at every pulse (these two facts are used in
Section~\ref{sec:elliptic-counting}). These are counts over \(\CC\).
\end{proposition}

In Table~\ref{tab:ks} the pulses are the four pairs with \(b=0\), and
their labels are the four that contain \(0\).

\begin{proof}
Let \(A\) be a solution with label \(L\). If \(A(0)=0\), then
\(\overline A(0)=0\), so \(0\in L\). Conversely, let \(0\in L\). At the
solution \(A\) we have \(\ell S_A=A\cdot\ell C_A\). By \textup{(V)} and
\eqref{eq:prime-collapse}, \(\ell C_A\) is \(\ell\)-integral and reduces
to \((-1)^d\overline C\), where
\(\overline C=(\rho^\ell-\rho)/\overline A\). Since \(\rho\) divides
\(\overline A\) and \(\rho^\ell-\rho\) has no repeated factor,
\(\overline C(0)\ne0\). The number \(\ell A(0)\) reduces to zero. Hence
\(\ell\bigl(A(0)+C_A(0)\bigr)\) reduces to \((-1)^d\overline C(0)\ne0\).
So \(A(0)+C_A(0)\ne0\), and \(R_0=0\) gives \(A(0)=0\), by
\eqref{eq:constant-remainder}: the pair is a pulse. The number of
\(d\)-element subsets of \(\FF_\ell\) that contain \(0\) is
\(\binom{\ell-1}{d-1}=S_p\).

At a pulse, \eqref{eq:constant-remainder} shows that the row of the
Jacobian matrix of the matching equations that belongs to \(R_0\) has a
single entry that can be different from zero, namely \(-C_A(0)\) in the
column of \(\alpha_d\). The Jacobian matrix is invertible, by
Theorem~\ref{thm:prime}(a). Hence \(C_A(0)\ne0\), so that
\(P_A(0)=s_pC_A(0)/2\ne0\), and the minor obtained by deleting that row
and that column is invertible. This minor is the Jacobian matrix of the
pulse equations. (For \(d=1\) there are no pulse equations, and only
\(C_A(0)\ne0\) is asserted.)
\end{proof}
